% =====================================================================
%  Appendix fragment -- \input from the thesis.
%  Uses only packages already loaded by the thesis:
%  amsmath, amssymb, mathtools, siunitx, hyperref, biblatex.
%  No booktabs, no empheq, no extra macros.
% =====================================================================

\section{Notation and conventions}
\label{app:notation}

\begin{center}
\begin{tabular}{ll}
\hline
$N_x,N_y$ & number of tones per axis, indices $i,j$\\
$W$ & frequency span per axis, tone spacing $\Delta f_{x}=W/(N_x-1)$\\
$\phi_{x,i},\phi_{y,j}$ & tone phases, free parameters at the AWG\\
$a_{x,i},a_{y,j}$ & tone \emph{power} weights, $A_s=\sqrt{a_{x,i}a_{y,j}}$\\
$u(r)$ & focal mode, normalised to $u(0)=1$\\
$f_0,\,T_0$ & fundamental beat frequency and period\\
$k_s$ & frequency order of spot $s$, $H_m$ group field of order $m$\\
$D_d$ & Fourier coefficient of the intensity, beat order $d$\\
$C$ & crest factor of the RF signal per axis\\
$\eta,\ \alpha,\ P_\mathrm{s}$ & diffraction efficiency, peak efficiency, saturation power\\
\hline
\end{tabular}
\end{center}

The amplitude ratios $r_x,r_y$ are \emph{power} ratios. The RF voltage
amplitudes therefore carry the square root,
\begin{equation}
a=[\,\varrho,1,\dots,1,\varrho\,]\ \text{in voltage},
\qquad
[\,r,1,\dots,1,r\,]\ \text{in power},
\qquad
\varrho:=\sqrt{r} .
\label{eq:app_conv}
\end{equation}

\section{Beating model}
\label{app:beating}

\begin{equation}
f_s=f_{x,i}+f_{y,j}=f_{\min}+k_sf_0,
\qquad
f_0=\frac{W}{\operatorname{lcm}(N_x-1,N_y-1)},
\qquad
T_0=\frac{1}{f_0} .
\end{equation}
\begin{equation}
g_s(\mathbf{r})=A_s\,u(\mathbf{r}-\mathbf{c}_s),
\qquad
\phi_s=\phi_{x,i}+\phi_{y,j},
\qquad
H_m(\mathbf{r})=\sum_{s:\,k_s=m}g_s(\mathbf{r})\,e^{i\phi_s} .
\end{equation}
Only $N_x+N_y$ of the $N_xN_y$ spot phases are free. Dropping the global factor
$e^{i2\pi f_{\min}t}$, which cancels in $|E|^2$,
\begin{equation}
\boxed{\;
E(\mathbf{r},t)=\sum_m H_m(\mathbf{r})\,e^{i2\pi mf_0t},
\qquad
I(\mathbf{r},t)=\sum_d D_d(\mathbf{r})\,e^{i2\pi df_0t}\;}
\label{eq:app_fourier}
\end{equation}
\begin{equation}
D_d(\mathbf{r})=\!\!\sum_{k_s-k_{s'}=d}\!\!
g_s(\mathbf{r})\,g_{s'}(\mathbf{r})\,e^{i(\phi_s-\phi_{s'})},
\qquad
D_{-d}=D_d^{*},
\qquad
f_\mathrm{beat}=d\,f_0 .
\end{equation}
Time average over one period $T_0$:
\begin{equation}
\langle I\rangle_t=D_0
=\sum_s g_s^2
+2\!\!\sum_{\substack{s<s'\\ k_s=k_{s'}}}\!\! g_s\,g_{s'}\cos(\phi_s-\phi_{s'}),
\label{eq:app_timeavg}
\end{equation}
so the time average equals the incoherent sum only if every group contains a
single spot; a degenerate pair contributes a static term that vanishes for
$\phi_s-\phi_{s'}=\pm\pi/2$. Different orders oscillate at different
frequencies and are orthogonal over $T_0$, so the variances add:
\begin{equation}
\operatorname{Var}_t[I]=\sum_{d\neq0}|D_d|^2
=\sum_{d>0}\sigma_d^2,
\qquad
\sigma_d^2:=2\,|D_d|^2 ,
\label{eq:app_parseval}
\end{equation}
\begin{equation}
U_t(\mathbf{r})=\frac{\sigma_I(\mathbf{r})}{\langle I\rangle_t(\mathbf{r})},
\qquad
\bar U_t=\big\langle U_t^2\big\rangle_R^{1/2}
=\Big(\sum_{d>0}\bar u_d^2\Big)^{1/2},
\qquad
\bar u_d=\bigg\langle\frac{\sigma_d^2}{\langle I\rangle_t^2}\bigg\rangle_R^{1/2},
\label{eq:app_Ut}
\end{equation}
with $\langle\,\cdot\,\rangle_R$ the average over the spot square $R$, the same
region as for $U_h$. The extreme order $d=k_{\max}-k_{\min}$ is produced by a
single pair of corner spots, so $|D_{d_{\max}}|$ cannot be changed by any
choice of tone phases.

\section{Choice of the tone phases}
\label{app:phases}

\subsection{The centring criterion}

The atom must see a pattern that is symmetric at \emph{every} instant, not
merely on average:
\begin{equation}
I(-\mathbf{r},t)=I(\mathbf{r},t)\quad\forall\,t
\qquad\Longleftrightarrow\qquad
E(-\mathbf{r},t)=e^{i\alpha(t)}E^{*}(\mathbf{r},t),
\label{eq:app_centred}
\end{equation}
with one and the same $\alpha(t)$ for the whole field. Starting from
\begin{equation}
E(\mathbf{r},t)=\sum_{n,m}a_nb_m\,u(\mathbf{r}-\mathbf{c}_{nm})\,
e^{-i2\pi(f_n+g_m)t}\,e^{i(\varphi_n+\psi_m)},
\end{equation}
three symmetries of the profile enter, with the mirror index $\bar n=N-1-n$:
\begin{align}
\mathbf{c}_{\bar n\bar m}&=-\mathbf{c}_{nm}
&&\text{(centred spot grid)},\label{eq:app_symgeo}\\
a_{\bar n}&=a_n,\qquad b_{\bar m}=b_m
&&\text{(amplitudes }[\varrho,1,\dots,1,\varrho]\text{)},\label{eq:app_symamp}\\
f_{\bar n}&=F-f_n,\qquad g_{\bar m}=G-g_m,
&&F=f_n+f_{\bar n}=2f_\mathrm{off}+(N-1)\Delta f .\label{eq:app_symfreq}
\end{align}
Evaluating $E$ at $-\mathbf{r}$ and relabelling the sum by the mirror index:

\medskip
\noindent\textbf{Step 1} (positions, $u$ even):
$u(-\mathbf{r}-\mathbf{c}_{\bar n\bar m})
=u\big(-(\mathbf{r}-\mathbf{c}_{nm})\big)=u(\mathbf{r}-\mathbf{c}_{nm})$,
the envelope factor is reproduced.

\noindent\textbf{Step 2} (amplitudes): $a_{\bar n}b_{\bar m}=a_nb_m$, unchanged.

\noindent\textbf{Step 3} (frequencies), the decisive one:
\begin{equation}
e^{-i2\pi(f_{\bar n}+g_{\bar m})t}
=e^{-i2\pi(F+G)t}\cdot e^{+i2\pi(f_n+g_m)t},
\label{eq:app_conj}
\end{equation}
i.e.\ the frequency mirror does not reproduce the time factor, it
\emph{conjugates} it, up to one global factor common to all spots.

\noindent\textbf{Step 4} (collect and compare):
\begin{align}
E(-\mathbf{r},t)&=e^{-i2\pi(F+G)t}\sum_{n,m}a_nb_m\,
u(\mathbf{r}-\mathbf{c}_{nm})\,e^{+i2\pi(f_n+g_m)t}\,
e^{i(\varphi_{\bar n}+\psi_{\bar m})},\\
E^{*}(\mathbf{r},t)&=\sum_{n,m}a_nb_m\,
u(\mathbf{r}-\mathbf{c}_{nm})\,e^{+i2\pi(f_n+g_m)t}\,
e^{-i(\varphi_{n}+\psi_{m})} .
\end{align}
Both sums run over the same basis functions with the same real coefficients and
differ only by a constant phase per term. They are proportional, and hence
Eq.~\eqref{eq:app_centred} holds, if and only if
$e^{i(\varphi_{\bar n}+\psi_{\bar m})}
=e^{i\Theta}e^{-i(\varphi_n+\psi_m)}$ for all $n,m$ with a \emph{single}
constant $\Theta$. Separating the two axes,
\begin{equation}
\boxed{\;
\varphi_n+\varphi_{N_x-1-n}=c_x\ \ \forall n,
\qquad
\psi_m+\psi_{N_y-1-m}=c_y\ \ \forall m \;}
\label{eq:app_centring}
\end{equation}
Mirror-symmetric phases, $\varphi_{\bar n}=\varphi_n$, are a different
condition and do not centre the pattern; they leave residual centroid offsets
of \SIrange{44}{84}{\nano \metre} against $\sigma=\SI{108}{\nano \metre}$,
whereas Eq.~\eqref{eq:app_centring} gives $\SI{0.00}{\nano \metre}$.

\subsection{Structure of the admissible family}

With $\varphi_n=c/2+\chi_n$, Eq.~\eqref{eq:app_centring} becomes
\begin{equation}
\chi_{\bar n}=-\chi_n
\qquad\Longrightarrow\qquad
\dim(\text{family})=\Big\lfloor\frac{N}{2}\Big\rfloor
\quad\text{instead of}\quad N-1 .
\label{eq:app_antisym}
\end{equation}
For odd $N$ the map $n\mapsto\bar n$ has a fixed point, so the middle tone is
paired with \emph{itself} and its phase is pinned:
\begin{equation}
2\varphi_{(N-1)/2}=c
\qquad\Longrightarrow\qquad
\varphi_{(N-1)/2}=\frac{c}{2}\ \ \text{or}\ \ \frac{c}{2}+\pi .
\label{eq:app_mid}
\end{equation}
A ramp $\chi_n=\delta\,\hat n$ with $\hat n:=n-(N-1)/2$ satisfies
Eq.~\eqref{eq:app_antisym} automatically and is a pure time shift of the
envelope,
\begin{equation}
A_\delta(t)=A_0(t+\tau),
\qquad
\tau=\frac{\delta}{2\pi\,\Delta f} ,
\label{eq:app_ramp}
\end{equation}
hence crest neutral and redundant with the trigger $t_0$. One of the
$\lfloor N/2\rfloor$ dimensions is therefore always inactive:
\begin{equation}
\dim(\text{crest-relevant})=\Big\lfloor\frac{N}{2}\Big\rfloor-1
\ \xrightarrow{\ N=3\ }\ 0 .
\label{eq:app_dim}
\end{equation}

\subsection{Crest factor}

\begin{equation}
C=\frac{\max_t|s(t)|}{\mathrm{rms}_t[s(t)]}
=\sqrt{2}\,\frac{\max_t|A(t)|}{\sqrt{\sum_n a_n^2}},
\qquad
\sqrt{2}\le C\le\sqrt{2}\,\frac{\sum_n a_n}{\sqrt{\sum_n a_n^2}}
\ \xrightarrow{a_n\equiv1}\ \sqrt{2N},
\label{eq:app_crest}
\end{equation}
both operations taken over a full envelope period; the denominator is
phase independent \cite{Schroeder1970,Boyd1986}. For $N_x=3$, inserting
Eq.~\eqref{eq:app_centring} as
$\varphi=\big(\tfrac{c}{2}-\chi,\ \tfrac{c}{2}\ \text{or}\ \tfrac{c}{2}+\pi,\
\tfrac{c}{2}+\chi\big)$ and factoring out $e^{ic/2}$, with
$\theta:=2\pi\Delta f_x t$,
\begin{equation}
\boxed{\;
A(t)\,e^{-ic/2}
=\varrho_x e^{-i(\chi+\theta)}\pm1+\varrho_x e^{+i(\chi+\theta)}
=2\varrho_x\cos(\theta+\chi)\pm1 \;}
\label{eq:app_n3}
\end{equation}
The only free parameter appears \emph{inside} the cosine, and $\cos$ sweeps
$[-1,1]$ over one period regardless of $\chi$, so
\begin{equation}
\max_t|A|=\max_{v\in[-1,1]}\big|2\varrho_x v\pm1\big|=2\varrho_x+1=\sum_n a_n
\qquad\text{for all }c,\chi\text{ and both branches},
\end{equation}
which is the upper bound of Eq.~\eqref{eq:app_crest}. With $r_x=0.97$ and
$\varrho_x=\sqrt{r_x}=0.9848858$,
\begin{equation}
C_x=\sqrt{2}\,\frac{2\varrho_x+1}{\sqrt{2r_x+1}}
=\sqrt{2}\,\frac{2.9697716}{1.7146428}=2.4494
\qquad\big(\sqrt{2N}=\sqrt{6}=2.4495\big).
\end{equation}
For $N_y=4$ there is no fixed point and Eq.~\eqref{eq:app_centring} leaves
$\psi=[0,b,c_y-b,c_y]$ with two parameters, one of them crest relevant. At
$c_y=114^\circ$:

\begin{center}
\begin{tabular}{cccl}
\hline
$b$ & $\psi$ [deg] & $\max|A|$ & $C_y$\\ \hline
$38^\circ$ & $[0,38,76,114]$ & $4.154$ & $2.8265$ \ (ramp, $=\sum_n a_n$)\\
$60^\circ$ & $[0,60,54,114]$ & $4.022$ & $2.7364$\\
$98^\circ$ & $[0,98,16,114]$ & $3.228$ & $2.1963$\\
$114^\circ$ & $[0,114,0,114]$ & $3.639$ & $2.4761$\\ \hline
\end{tabular}
\end{center}

A power amplifier clips the instantaneous voltage and reaches compression long
before any thermal limit \cite{Pozar2012,Cripps2006}, so with
$P_\mathrm{peak}=C^2P_\mathrm{avg}$ and
$P_\mathrm{peak}\le P_\mathrm{1dB}$,
\begin{equation}
\boxed{\;
P_\mathrm{avg}\le\frac{P_\mathrm{1dB}}{C^{2}}\;}
\qquad
P_\mathrm{1dB}=\SI{36}{\decibelmilliwatt}=\SI{3.98}{\watt}
\ \text{\cite{MiniCircuitsZHL}} .
\label{eq:app_backoff}
\end{equation}

\begin{center}
\begin{tabular}{lccc}
\hline
& $C$ & $C^2$ & $P_\mathrm{avg}\le P_\mathrm{1dB}/C^2$\\ \hline
$x$ axis, $N_x=3$, forced by Eq.~\eqref{eq:app_dim}
 & $2.4494$ & $6.00$ & \SI{0.66}{\watt}\\
$y$ axis, $N_y=4$, $b=98^\circ$, minimum in the family
 & $2.1963$ & $4.82$ & \SI{0.83}{\watt}\\ \hline
\end{tabular}
\end{center}

\section{Pulse area for a given profile power}
\label{app:pulsearea}

The optical power $P$ in the profile fixes the scale $s$ in
$I=s\tilde I$, where $\tilde I=|E|^2$ is the simulated intensity in units of
$u(0)=1$:
\begin{equation}
P=\int\langle I\rangle_t\,\mathrm{d}^2r=s\,\tilde P_\mathrm{prof},
\qquad
\tilde P_\mathrm{prof}=\sum_s a_sP_1
+2\!\!\sum_{\substack{s<s'\\ k_s=k_{s'}}}\!\!A_sA_{s'}
\cos(\phi_s-\phi_{s'})\,O(\varrho_{ss'}),
\label{eq:app_scale}
\end{equation}
with the single-spot power $P_1=\int u^2\,\mathrm{d}^2r$ and the mode
autocorrelation $O$ evaluated at the spot separation $\varrho_{ss'}$
\cite{Goodman2017,BornWolf1999,DLMF}:
\begin{equation}
P_1^\mathrm{G}=\frac{\pi w_0^2}{2},
\quad
O^\mathrm{G}(\varrho)=P_1^\mathrm{G}\,e^{-\varrho^2/2w_0^2};
\qquad
P_1^\mathrm{A}=\frac{4\pi}{k^2},
\quad
O^\mathrm{A}(\varrho)=P_1^\mathrm{A}\,u(\varrho),
\label{eq:app_overlap}
\end{equation}
$k=3.8317/(1.4830\,w_0)$. The Airy autocorrelation is the mode itself, sign
included. Adiabatic elimination of the $5P$ manifold, with signed dipole
matrix elements and hyperfine shifts from ARC \cite{ARC2017,SteckRb85} and the
far detuned scaling of \cite{Grimm2000}, makes all three atomic quantities
linear in intensity:
\begin{equation}
\Omega=C_\mathrm{rabi}I,
\qquad
\delta_\mathrm{LS}=C_\mathrm{shift}I,
\qquad
\Gamma_\mathrm{sc}=C_\mathrm{scatter}I,
\qquad
\varepsilon_\mathrm{LS}=\frac{\delta_\mathrm{LS}}{\Omega}=\text{const} .
\end{equation}
The area of a rectangular pulse on $[t_0,t_0+T_p]$ is therefore
\begin{equation}
\boxed{\;
\theta(\mathbf{r},t_0)=C_\mathrm{rabi}\,P\,
\frac{\displaystyle\int_{t_0}^{t_0+T_p}\tilde I(\mathbf{r},t)\,\mathrm{d}t}
     {\displaystyle\int\langle\tilde I\rangle_t\,\mathrm{d}^2r}\;}
\label{eq:app_area}
\end{equation}
in which the arbitrary mode normalisation cancels. Inserting
Eq.~\eqref{eq:app_fourier} and integrating each order over the window,
\begin{equation}
\int_{t_0}^{t_0+T_p}\!\!e^{i2\pi df_0t}\,\mathrm{d}t
=T_p\,\operatorname{sinc}(d f_0T_p)\,e^{i2\pi df_0(t_0+T_p/2)},
\label{eq:app_sinc}
\end{equation}
\begin{equation}
\theta(\mathbf{r},t_0)=\frac{C_\mathrm{rabi}PT_p}{\tilde P_\mathrm{prof}}
\bigg[D_0+2\sum_{d>0}|D_d|\,\operatorname{sinc}(df_0T_p)\,
\cos\!\Big(2\pi df_0\big(t_0+\tfrac{T_p}{2}\big)+\arg D_d\Big)\bigg].
\label{eq:app_areaclosed}
\end{equation}
For $T_p\gg T_0$ every $d\neq0$ is damped and the pulse sees the time average;
for $T_p\ll T_0$ the sinc is unity and the pulse samples the beating, so
$\theta$ depends on the trigger $t_0$. Averaging over the atom, with
$W(\mathbf{r})$ the thermal position distribution of width
$\sigma^2=\frac{\hbar}{2m\omega_r}\coth\frac{\hbar\omega_r}{2k_BT}$
\cite{Tuchendler2008},
\begin{equation}
c_d=\int W\,D_d\,\mathrm{d}^2r,
\qquad
U(\theta)=\frac{\sigma_W(\theta)}{\langle\theta\rangle_W},
\qquad
\langle p_\uparrow\rangle_W=\int W\sin^2\frac{\theta}{2}\,\mathrm{d}^2r
\ \neq\ \sin^2\frac{\langle\theta\rangle_W}{2} .
\label{eq:app_atom}
\end{equation}
Since $\theta$ is linear in the $D_d$, a full scan over $t_0$ is a single
matrix-vector product in the $c_d$.

\begin{center}
\begin{tabular}{lr@{\qquad}lr}
\hline
$\tilde P_\mathrm{prof}$ & \SI{2.615e-11}{\metre\squared}
 & $\langle\theta\rangle_W$ at best trigger & $11.54\,\pi$\\
$C_\mathrm{rabi}$ & \SI{11.66}{\radian\per\second}/(W/m$^2$)
 & $U(\theta)$ there & \SI{0.72}{\percent}\\
$\varepsilon_\mathrm{LS}$ & $-0.061$
 & $P$ for a \SI{1}{\micro \second} $\pi$ pulse & \SI{0.87}{\micro \watt}\\
\hline
\end{tabular}
\end{center}

\noindent Working point: $3\times4$, Airy mode, $w_0=\SI{1.04}{\micro \metre}$,
$r_x/r_y=0.97/1.16$, $\Delta=+\SI{50}{\giga \hertz}$,
$P=\SI{10}{\micro \watt}$ in the profile, $T_p=\SI{1}{\micro \second}$,
$\sigma=\SI{108}{\nano \metre}$.

\section{Diffraction efficiency and profile power}
\label{app:eta}

Single grating in the Bragg regime, with
$\Delta n\propto\sqrt{P_\mathrm{ac}}$ from the photoelastic effect
\cite{Korpel1981,SalehTeich2019,AAOptoTheory}:
\begin{equation}
\eta(P)=\alpha(\nu)\,
\sin^2\!\bigg(\frac{\pi}{2}\sqrt{\frac{P}{P_\mathrm{s}(\nu)}}\bigg),
\qquad
\kappa=\frac{\pi}{\lambda}\sqrt{\frac{M_2L}{2H}}\sqrt{\zeta},
\qquad
P_\mathrm{s}=\Big(\frac{\pi}{2\kappa}\Big)^{2} .
\label{eq:app_etasingle}
\end{equation}
With $N$ tones the gratings share one undiffracted wave; the coupled wave
equations \cite{Kogelnik1969,Laude2003} reduce exactly to a two wave problem
and give \cite{msc:mittenbuehler,Antonov2007}
\begin{equation}
\boxed{\;
\eta_\mathrm{tot}=\alpha\,
\sin^2\!\bigg(\frac{\pi}{2}\sqrt{\frac{P_\mathrm{tot}}{P_\mathrm{s}}}\bigg),
\qquad
\eta_n=\frac{P_n}{P_\mathrm{tot}}\,\eta_\mathrm{tot},
\qquad
P_\mathrm{tot}=\sum_n P_n \;}
\label{eq:app_multitone}
\end{equation}
so the total efficiency depends only on the total RF power and not on its
distribution among the tones, and the tone phases do not enter at all.
Combining with Eq.~\eqref{eq:app_backoff},
\begin{equation}
\boxed{\;
\eta_\mathrm{axis}(C)=\alpha\,
\sin^2\!\bigg(\frac{\pi}{2C}\sqrt{\frac{P_\mathrm{1dB}}{P_\mathrm{s}}}\bigg),
\qquad
\eta_\mathrm{profile}=\eta_x\,\eta_y\,\Xi \;}
\label{eq:app_etacrest}
\end{equation}
with the measured cross axis term
\begin{equation}
\Xi_{ij}=\Big(1-\xi_0\big(\nu^{(x)}_i-\xi_x\big)^2
\big(\nu^{(y)}_j-\xi_y\big)^2\Big)^{2},
\qquad
\xi_x=\SI{108.0}{\mega \hertz},
\quad
\xi_y=\SI{102.3}{\mega \hertz},
\label{eq:app_xi}
\end{equation}
which falls to $0.80$ in the corners of the full scan range but is common to
all twelve spots over the \SI{0.37}{\mega \hertz} span used here
\cite{msc:mittenbuehler}. With the parameters measured on this device family,
$\alpha\approx0.92$ per axis and
$P_\mathrm{s}\approx\SI{1.70}{\watt}$ resp.\ $\SI{1.62}{\watt}$ near
\SI{100}{\mega \hertz} \cite{msc:mittenbuehler}:

\begin{center}
\begin{tabular}{lccc}
\hline
& $C$ & $P_\mathrm{avg}$ & $\eta_\mathrm{axis}$\\ \hline
$x$ axis, 3 tones, centred, forced & $2.449$ & \SI{0.66}{\watt} & $0.636$\\
$y$ axis, 4 tones, $b=98^\circ$ & $2.196$ & \SI{0.83}{\watt} & $0.746$\\
$y$ axis, ramp, worst in the family & $2.827$ & \SI{0.50}{\watt} & $0.538$\\
hypothetical $C=\sqrt{2}$ & $1.414$ & \SI{1.99}{\watt} & $0.905$\\ \hline
\multicolumn{3}{l}{$\eta_\mathrm{profile}=\eta_x\eta_y$ at the working point}
 & \SI{47.4}{\percent}\\
\multicolumn{3}{l}{at $C=\sqrt{2}$ on both axes, a factor $1.71$}
 & \SI{80.9}{\percent}\\
\multicolumn{3}{l}{laser power for \SI{10}{\micro \watt} in the profile}
 & \SI{21.1}{\micro \watt}\\
\multicolumn{3}{l}{for the \SI{0.87}{\micro \watt} of a $\pi$ pulse}
 & \SI{1.8}{\micro \watt}\\ \hline
\end{tabular}
\end{center}

\noindent Assumptions: common Bragg matching for all tones, the
\SI{0.37}{\mega \hertz} span costing \SI{0.007}{\percent} against the
\SI{36}{\mega \hertz} bandwidth by the detuned two wave solution
\cite{Kogelnik1969}; first order only; a linear transducer, the instantaneous
strain carrying the crest factor as well, so that a high $C$ also brings on
acoustic nonlinearity; isotropic diffraction, whereas the device is
anisotropic and birefringent, the two agreeing qualitatively \cite{Chang1995};
and a two tone estimate for the intermodulation, which is optimistic once many
tones contribute to the same bin.
